\documentclass[10pt]{scrartcl}
\usepackage[mdthm,colorsec,noauthor,noasy]{evan}
\usepackage{handout}
\newcommand{\aug}[2]{\left[\begin{array}{@{}#1@{}}#2\end{array}\right]}
\newcommand{\vv}[1]{\mathbf{#1}}
\DeclareMathOperator{\spn}{Span}\DeclareMathOperator{\Nul}{Nul}\DeclareMathOperator{\Col}{Col}\DeclareMathOperator{\Row}{Row}\DeclareMathOperator{\tr}{tr}
\title{Math 54 --- Week 8}
\subtitle{Solutions to the problem set}
\author{}\date{}\ihead{}
\begin{document}
\maketitle

\begin{prob}\textbf{\xlink{math54}{2026-10-12-diagonalization-complex-eigenvalues}{prob-P1}{P1}.}\par\idea{Find eigenvalues and eigenvectors, then use $A^k=PD^kP^{-1}$.}
\begin{steps}
  \ii $\tr A=3$ and $\det A=2-12=-10$, so $\det(A-\lambda I)=\lambda^2-3\lambda-10=(\lambda-5)(\lambda+2)$, giving eigenvalues $5$ and $-2$.
  \ii For $\lambda=5$, Row~$1$ of $A-5I$ is $(-3,3) \implies \vv p_1=(1,1)$. For $\lambda=-2$, Row~$1$ of $A+2I$ is $(4,3) \implies \vv p_2=(-3,4)$.
  \ii Since $\det P=4+3=7$:
  \[
    P=\begin{bmatrix}1&-3\\1&4\end{bmatrix},\qquad
    D=\begin{bmatrix}5&0\\0&-2\end{bmatrix},\qquad
    P^{-1}=\frac{1}{7}\begin{bmatrix}4&3\\-1&1\end{bmatrix}.
  \]
  \ii Therefore,
  \[
    A^k = PD^kP^{-1} = \frac{1}{7}\begin{bmatrix}5^k&-3(-2)^k\\5^k&4(-2)^k\end{bmatrix}\begin{bmatrix}4&3\\-1&1\end{bmatrix}
    = \frac{1}{7}\begin{bmatrix}4\cdot 5^k+3(-2)^k & 3\cdot 5^k-3(-2)^k\\ 4\cdot 5^k-4(-2)^k & 3\cdot 5^k+4(-2)^k\end{bmatrix}.
  \]
  (At $k=1$ this recovers $\frac{1}{7}\begin{psmallmatrix}14&21\\28&7\end{psmallmatrix}=A$.)
\end{steps}
\end{prob}

\begin{prob}\textbf{\xlink{math54}{2026-10-12-diagonalization-complex-eigenvalues}{prob-P2}{P2}.}\par\idea{Compare the eigenspace dimension against the multiplicity for each repeated eigenvalue (\xlink{math54}{2026-10-12-diagonalization-complex-eigenvalues}{thm-lay-5-7}{Thm.~5.7}).}
\begin{steps}
  \ii \textbf{For $A$:} Lower triangular with eigenvalues $2,2,3$. Here
  \[
    A-2I=\begin{bmatrix}0&0&0\\1&0&0\\0&0&1\end{bmatrix}
  \]
  has rank~$2$, so $\dim\Nul(A-2I)=3-2=1<2$. Thus $A$ is \textbf{not diagonalizable}.
  \ii \textbf{For $B$:} Expanding $\det(B-\lambda I)$ down column~$2$ gives $(5-\lambda)(4-\lambda)(5-\lambda)=(5-\lambda)^2(4-\lambda)$, so eigenvalues are $5,5,4$.
  For $\lambda=5$: $B-5I$ has rows $(-1,0,-2)$, $(2,0,4)$, $(0,0,0)$, rank~$1$, $\dim E(5)=2$ with basis $(0,1,0),(-2,0,1)$.
  For $\lambda=4$: $B-4I$ gives $x_3=0$ and $2x_1+x_2=0 \implies (-1,2,0)$.
  Hence $B=PDP^{-1}$ with
  \[
    P=\begin{bmatrix}-1&0&-2\\2&1&0\\0&0&1\end{bmatrix},\qquad
    D=\begin{bmatrix}4&0&0\\0&5&0\\0&0&5\end{bmatrix}.
  \]
\end{steps}
\end{prob}

\begin{prob}\textbf{\xlink{math54}{2026-10-12-diagonalization-complex-eigenvalues}{prob-P3}{P3}.}\par\idea{Sum the eigenspace dimensions (\xlink{math54}{2026-10-12-diagonalization-complex-eigenvalues}{thm-lay-5-7}{Thm.~5.7}).}
Eigenvalues $2$ and $3$ have multiplicity~$1$, so each eigenspace has dimension~$1$ (Thm.~7a). If $\rank(A-I)=3$, then $\dim E(1)=4-3=1$, total $1+1+1=3<4$ (\textbf{not diagonalizable}). If $\rank(A-I)=2$, then $\dim E(1)=2$, total $2+1+1=4$ (\textbf{diagonalizable}).
\end{prob}

\begin{prob}\textbf{\xlink{math54}{2026-10-12-diagonalization-complex-eigenvalues}{prob-P4}{P4}.}\par\idea{Compute $[T(1)]_{\mathcal C}$ and $[T(t)]_{\mathcal C}$ as columns.}
$T(1)=t+1\leftrightarrow (1,1,0)^{\mathsf T}$ and $T(t)=t+t^2\leftrightarrow (0,1,1)^{\mathsf T}$, so
\[
  M = \begin{bmatrix}1&0\\1&1\\0&1\end{bmatrix},\qquad
  [T(3-2t)]_{\mathcal C} = M\begin{bmatrix}3\\-2\end{bmatrix} = \begin{bmatrix}3\\1\\-2\end{bmatrix}
  \implies T(3-2t)=3+t-2t^2.
\]
(Direct check: $(t+1)(3-2t)=3+t-2t^2$.)
\end{prob}

\begin{prob}\textbf{\xlink{math54}{2026-10-12-diagonalization-complex-eigenvalues}{prob-P5}{P5}.}\par\idea{Compute $A\vv b_1$ and $A\vv b_2$ and express them in $\mathcal B$-coordinates (\xlink{math54}{2026-10-12-diagonalization-complex-eigenvalues}{thm-lay-5-8}{Thm.~5.8}).}
Notice $A(1,1)=(2,2)=2(1,1)$ and $A(1,2)=(1,2)=1\cdot(1,2)$. Their $\mathcal B$-coordinates are $(2,0)^{\mathsf T}$ and $(0,1)^{\mathsf T}$, so the $\mathcal B$-matrix is $\operatorname{diag}(2,1)$. Because it is diagonal, $(1,1)$ and $(1,2)$ are eigenvectors of $A$ with eigenvalues $2$ and $1$.
\end{prob}

\begin{prob}\textbf{\xlink{math54}{2026-10-12-diagonalization-complex-eigenvalues}{prob-P6}{P6}.}\par\idea{Find $\lambda=a-bi$, compute one eigenvector $\vv v\in\CC^2$, and apply \xlink{math54}{2026-10-12-diagonalization-complex-eigenvalues}{thm-lay-5-9}{Thm.~5.9}.}
\begin{steps}
  \ii $\tr A=6$ and $\det A=8+2=10 \implies \lambda^2-6\lambda+10=0 \implies \lambda=3\pm i$.
  \ii For $\lambda=3-i$, Row~$1$ of $A-(3-i)I$ is $(-1+i,\ -1)$, giving eigenvector $\vv v=(-1,1-i)$ (check: $A\vv v=(-3+i,\ 2-4i)=(3-i)\vv v$).
  \ii Taking $\operatorname{Re}\vv v=(-1,1)$ and $\operatorname{Im}\vv v=(0,-1)$ with $a=3,b=1$:
  \[
    P=\begin{bmatrix}-1&0\\1&-1\end{bmatrix},\qquad
    C=\begin{bmatrix}3&-1\\1&3\end{bmatrix},\qquad
    A=PCP^{-1}.
  \]
  \ii Scale factor: $r=\sqrt{3^2+1^2}=\sqrt{10}$; rotation angle: $\cos\varphi=3/\sqrt{10},\ \sin\varphi=1/\sqrt{10} \implies \varphi\approx 18.4^\circ$.
\end{steps}
\end{prob}

\begin{prob}\textbf{\xlink{math54}{2026-10-12-diagonalization-complex-eigenvalues}{prob-P7}{P7}.}\par\idea{Take inverses and transposes of $A=PDP^{-1}$.}
\begin{steps}
  \ii Since $A=PDP^{-1}$ is invertible, $0\neq\det A=\det P\det D\det P^{-1}=\det D=d_1\cdots d_n$, so every $d_i\neq 0$ and $D^{-1}=\operatorname{diag}(1/d_1,\dots,1/d_n)$ exists. Reversing the product order gives $A^{-1}=(P^{-1})^{-1}D^{-1}P^{-1}=PD^{-1}P^{-1}$, so $A^{-1}$ is diagonalizable (with the same $P$).
  \ii Taking transposes reverses order: $A^{\mathsf T}=(P^{-1})^{\mathsf T}D^{\mathsf T}P^{\mathsf T}=QDQ^{-1}$ where $Q=(P^{-1})^{\mathsf T}$ and $Q^{-1}=P^{\mathsf T}$, because $D^{\mathsf T}=D$. Thus $A^{\mathsf T}$ is \textbf{always} diagonalizable.
\end{steps}
\end{prob}

\begin{prob}\textbf{\xlink{math54}{2026-10-12-diagonalization-complex-eigenvalues}{prob-P8}{P8}.}\par\idea{Expand $\vv x_0$ in eigenvectors and let $k\to\infty$.}
Row sums equal $1 \implies \lambda_1=1$ with eigenvector $(1,1)$; also $A(-1,1)=\frac12(-1,1)\implies\lambda_2=\frac12$. Since $\vv x_0=(1,0)=\frac12(1,1)-\frac12(-1,1)$,
\[
  \vv x_k = A^k\vv x_0 = \frac12\begin{bmatrix}1\\1\end{bmatrix} - \frac12\left(\frac12\right)^k\begin{bmatrix}-1\\1\end{bmatrix} = \begin{bmatrix}\frac12+\frac{1}{2^{k+1}}\\[3pt]\frac12-\frac{1}{2^{k+1}}\end{bmatrix} \xrightarrow{k\to\infty} \begin{bmatrix}1/2\\1/2\end{bmatrix}.
\]
\end{prob}

\end{document}
